% Kandydat restrukturyzacji; nie jest aktywnym rozdziałem Atlasu.
\section{Transformacja Fouriera w przestrzeni $L^2$ i twierdzenie Plancherela}
\providecommand{\atlasref}[2]{#2}

\begin{twierdzenie}[Twierdzenie Plancherela]
\label{thm:plancherel22}
Transformacja Fouriera z normalizacją $(2\pi)^{-d/2}$
przedłuża się jednoznacznie do operatora unitarnego
$\mathcal F:L^2(\R^d;\C)\to L^2(\R^d;\C)$,
czyli dla $f,g\in L^2$ zachodzi
$\int_{\R^d}\widehat f(\xi)\overline{\widehat g(\xi)}
\,\mathrm d\lambda_d(\xi)
=\int_{\R^d}f(x)\overline{g(x)}\,\mathrm d\lambda_d(x)$.
Dla $f\in L^2$ zachodzi ponadto
$\mathcal F T_f=T_{\widehat f}$, gdzie $\widehat f$
oznacza transformatę w $L^2$.
\end{twierdzenie}
\begin{proof}
Najpierw ustalmy $f,\varphi\in\mathcal S$.
Ponieważ $f,\varphi\in L^1$, moduł funkcji
podcałkowej w całce podwójnej jest całkowalny:
\[
 \int_{\R^d}\int_{\R^d}|f(x)\varphi(\xi)|
 \,\mathrm d\lambda_d(x)\,\mathrm d\lambda_d(\xi)
 =\|f\|_1\|\varphi\|_1.
\]
Twierdzenie Fubiniego i symetria jądra
$e^{-ix\cdot\xi}$ dają
\[
 \int_{\R^d}\widehat f(\xi)\varphi(\xi)
 \,\mathrm d\lambda_d(\xi)
 =\int_{\R^d}f(x)(\mathcal F\varphi)(x)
 \,\mathrm d\lambda_d(x).
\]
Jest to zgodność $\mathcal F T_f=T_{\widehat f}$ na
$\mathcal S$; parowanie jest biliniowe, a nie
półtoraliniowe.

Niech teraz $f,g\in\mathcal S$. Stosujemy powyższą
tożsamość do $\varphi=\overline{\widehat g}\in\mathcal S$.
Ze wzoru odwrotnego w Twierdzeniu~\ref{thm:fourier-schwartz17}
wynika
$\mathcal F(\overline{\widehat g})(x)
=\overline{(\mathcal F^{-1}\widehat g)(x)}
=\overline{g(x)}$.
Otrzymujemy
\[
 \int_{\R^d}\widehat f(\xi)\overline{\widehat g(\xi)}
 \,\mathrm d\lambda_d(\xi)
 =\int_{\R^d}f(x)\overline{g(x)}\,\mathrm d\lambda_d(x).
\]
W szczególności $\|\widehat f\|_2=\|f\|_2$.
Przestrzeń $\mathcal S$ jest gęsta w $L^2$, bo zawiera
gęstą podprzestrzeń $\mathcal D$. Dla $f\in L^2$
wybieramy $f_n\in\mathcal S$ z $f_n\to f$ w $L^2$.
Ciąg $(\widehat f_n)$ jest Cauchy'ego w $L^2$,
ponieważ $\|\widehat f_n-\widehat f_m\|_2
=\|f_n-f_m\|_2$. Z zupełności istnieje granica;
dla innego ciągu przybliżającego $(g_n)$ tę samą granicę
zapewnia $\|\widehat f_n-\widehat g_n\|_2
=\|f_n-g_n\|_2\to0$. Definiujemy $\widehat f$ jako
tak uzyskaną granicę. Liniowość i zachowanie normy
przechodzą na granicę; jednoznaczność przedłużenia
wynika z gęstości $\mathcal S$.
Obraz izometrii jest domknięty: jeśli ciąg obrazów
jest zbieżny, ciąg przeciwobrazów jest Cauchy'ego i
jego granica ma żądany obraz. Obraz zawiera
$\mathcal F(\mathcal S)=\mathcal S$, więc jest całym
$L^2$. Zachowanie iloczynu skalarnego wynika z ciągłości
iloczynu skalarnego na $L^2$ i tożsamości na $\mathcal S$.
Otrzymujemy zatem operator unitarny.

Na koniec, dla $f_n\to f$ w $L^2$ oraz ograniczonego
$B\subseteq\mathcal S$ mamy
\[
 P_B(T_{f_n}-T_f)
 \le C_d\|f_n-f\|_2
 \sup_{\varphi\in B}Q_{d+1,0}(\varphi)\longrightarrow0.
\]
Analogicznie $T_{\widehat f_n}\to T_{\widehat f}$ silnie
w $\mathcal S'$, bo $\widehat f_n\to\widehat f$ w $L^2$.
Z ciągłości $\mathcal F$ w topologii silnej i
wykazanej zgodności na $\mathcal S$ otrzymujemy
$\mathcal FT_f=\lim_n\mathcal FT_{f_n}
=\lim_nT_{\widehat f_n}=T_{\widehat f}$.
Jednoznaczność granicy wynika z warunku Hausdorffa
dla topologii silnej na $\mathcal S'$.
\end{proof}

\begin{stwierdzenie}[Zgodność transformat całkowej, dystrybucyjnej i unitarnej]
\label{prop:dt-fourier-l1-l2-compat}
Dla $f\in L^1$ transformacja dystrybucyjna spełnia
$\mathcal FT_f=T_{\widehat f}$, gdzie $\widehat f$
jest transformatą całkową. Gdy $f\in L^1\cap L^2$,
transformata całkowa jest reprezentantem transformaty
unitarnej z Twierdzenia~\ref{thm:plancherel22}.
\end{stwierdzenie}
\begin{proof}
Dla $\varphi\in\mathcal S$ całkowalność modułu
$|f(x)\varphi(\xi)|$ wynika z $f,\varphi\in L^1$.
Z twierdzenia Fubiniego otrzymujemy
\[
 \int_{\R^d}\widehat f(\xi)\varphi(\xi)\,\mathrm d\lambda_d(\xi)
 =\int_{\R^d}f(x)(\mathcal F\varphi)(x)\,\mathrm d\lambda_d(x).
\]
Jest to definicyjna tożsamość transformacji dystrybucyjnej.
Dla $f\in L^1\cap L^2$ twierdzenie Plancherela daje
tę samą dystrybucję regularną dla transformaty w $L^2$.
Oba reprezentanty są lokalnie całkowalne; ich równość
prawie wszędzie wynika ze
Stwierdzenia~\ref{prop:regular-injective19}.
\end{proof}
